\section{Finite-window supersolutions with bounded increments}
\label{sc:section}

Let $P_t$ denote the inert-boundary semigroup.  An ordinary vertex
addition means a single leaf birth, which changes the core count by
either zero or one.
Write $G(t)$ for the supremum of newborn retained means over finite
environments and $H(t)=1+\int_0^tG(r)\,dr$.  By Section~\ref{uniform:section},
\[
 H(t+u)\le H(t)H(u),\qquad
 \lambda=\lim_{t\to\infty}\frac{\log H(t)}t,
 \qquad \mathbb E_y^\infty C(t)\le C(y)H(t).
\]
Each $H(T)$ for rational $T>0$ is effectively computable by the certified
finite-profile and finite-population reductions.

\begin{lemma}[Finite-time sensitivity]
\label{sc:sensitivity}
Let $\psi$ be any function with ordinary-birth increments bounded in
absolute value by $J$.  Assuming the expectations below exist, for every
ordinary vertex addition $y\mapsto y'$ and every $t\ge0$,
\begin{equation}\label{sc:sensitivity-bound}
 |P_t\psi(y')-P_t\psi(y)|\le J e^{2t}.
\end{equation}
\end{lemma}
\begin{proof}
Couple the two full vertex processes using common rate-one clocks at
common vertices and maximal couplings of their neighbour choices.
Initially all vertices of $y$ are common, and the added vertex is
unmatched.  Match newly born vertices only when their source event and
target agree in both copies; otherwise use distinct new labels.
The common vertices consequently induce a connected tree containing the
whole initial $y$.  Each component of unmatched vertices attaches to
this common tree through exactly one edge.

Let $Z(t)$ be the total number of unmatched vertices across both copies.
An unmatched vertex clock rings at rate one and creates at most one
new unmatched vertex.  At a common vertex, the probability of a
neighbour-choice mismatch is
\[
 \frac{\max(a,b)}{d_0+\max(a,b)}\le\frac{a+b}{2},
\]
where $d_0\ge1$ is the number of shared neighbours (counting the shared
inert stub at the root) and $a,b$ are the two counts of unmatched
neighbours.  The inequality follows directly for the integer cases
$a+b=1$ and $a+b\ge2$, and is trivial when $a+b=0$.
Summing over common vertices bounds the mismatch rate by half the
number of common--unmatched edges, and hence by $Z(t)/2$.
Each such mismatch creates at most two unmatched vertices.  Inert-stub
selections can only reduce these increments.  Hence
\[
 \frac{d}{dt}\mathbb EZ(t)\le2\mathbb EZ(t),
 \qquad Z(0)=1,
\]
so $\mathbb EZ(t)\le e^{2t}$.

At any time the two trees can each be built from their common connected
tree by ordinary leaf additions, with total number of additions $Z(t)$.
The common tree still contains the initial rooted core and inert stub,
so all those intermediate states are legal.  Telescoping the jump bound
gives $|\psi(X_t')-\psi(X_t)|\le JZ(t)$.  Taking expectations proves
\eqref{sc:sensitivity-bound}.
\end{proof}

\begin{theorem}[A finite time-window supersolution]
\label{sc:window}
Let $s,T>0$, and suppose a certified rational $q<1$ satisfies
\begin{equation}\label{sc:time-test}
 e^{-sT}H(T)\le q.
\end{equation}
Define
\begin{equation}\label{sc:window-U}
 \boxed{U_{s,T}(y)=\frac1{1-q}
                \int_0^T e^{-st}\mathbb E_y^\infty C(t)\,dt.}
\end{equation}
Then $U_{s,T}$ is nonnegative and, for computable $s,T$, effectively
computable at every finite state.  It satisfies
\begin{equation}\label{sc:window-properties}
 (s-\mathcal L_\infty)U_{s,T}\ge C,
 \qquad
 |U_{s,T}(y')-U_{s,T}(y)|
 \le B_{s,T}:=\frac1{1-q}\int_0^T e^{(2-s)t}\,dt
\end{equation}
for every ordinary vertex addition.  The final integral equals
$(e^{(2-s)T}-1)/(2-s)$ when $s\ne2$, and $T$ when $s=2$.
Moreover
\begin{equation}\label{sc:window-linear}
 0\le U_{s,T}(y)\le a_{s,T}C(y),\qquad
 a_{s,T}:=\frac1{1-q}\int_0^T e^{-st}H(t)\,dt.
\end{equation}
The elementary bound $H(t)\le e^t$ permits the explicit replacement
$a_{s,T}\le\int_0^T e^{(1-s)t}dt/(1-q)$ whenever only an upper
constant is needed.
\end{theorem}
\begin{proof}
The finite-horizon resolvent identity gives
\[
 (s-\mathcal L_\infty)
      \int_0^T e^{-st}P_tC\,dt
 =C-e^{-sT}P_TC.
\]
It is justified pointwise by finite state localization and Yule
domination.  The conditional old-core mean bound gives
$P_TC(y)\le C(y)H(T)$.
The right side of the resolvent identity is consequently at least
$(1-q)C$, proving the supersolution inequality.
Lemma~\ref{sc:sensitivity} applied to $\psi=C$, $J=1$ proves the
increment estimate after integration.  The old-core mean bound also
gives \eqref{sc:window-linear}.

At a fixed finite initial state, rational finite-state generator
truncations compute $P_tC$ uniformly over $[0,T]$: the real population
is dominated by a Yule process, and $C\le N$ gives an explicit
discarded-population error.  Uniformization supplies rational intervals
for the matrix-exponential entries.  Time quadrature is effective as
well, since $0\le\mathcal L_\infty C\le N$ and
$\mathbb E N(t)\le N(0)e^t$ bound its finite-horizon derivative.
Thus the integral defining $U_{s,T}$ is computable with certified error.
\end{proof}

\begin{corollary}[Completeness above $\lambda$]
\label{sc:complete-lambda}
For every rational $s>\lambda$, a search over positive integer $T$ and
certified upper approximations of $H(T)$ eventually verifies
\eqref{sc:time-test} with $q=1/2$, and hence constructs a supersolution
with a uniform
ordinary-birth increment at $s$.
\end{corollary}
\begin{proof}
Since $s>\lambda$, $e^{-sT}H(T)\to0$.  Thus the strict inequality in
\eqref{sc:time-test} holds for sufficiently large integer $T$ with
$q=1/2$.  Dovetailing the positive integers $T$ and approximation
precisions eventually certifies that strict inequality.  The preceding
theorem then supplies the explicit function and increment bound.
\end{proof}

